\section{Discussion}

In this work we closely analyze hallucination in object captioning models.
Our work is similar to other works which attempt to characterize flaws of different evaluation metrics~\cite{kilickaya2016re}, though we focus specifically on hallucination.
Likewise, our work is related to other work which aims to build better evaluation tools (\cite{cider}, ~\cite{spice}, ~\cite{cui2018learning}).
However, we focus on carefully quantifying and characterizing one important type of error: object hallucination.

A significant number of objects are hallucinated in current captioning models (between 5.5\% and 13.1\% of MSCOCO objects).
Furthermore, hallucination does not always agree with the output of standard captioning metrics. 
For instance, the popular self critical loss increases CIDEr score, but also the amount of hallucination.
Additionally, we find that given two sentences with similar CIDEr, SPICE, or METEOR scores from two different models, the number of hallucinated objects might be quite different.
This is especially apparent when standard metrics assign a low score to a generated sentence.
Thus, for challenging caption tasks on which standard metrics are currently poor (e.g., the LSMDC dataset~\cite{rohrbach2017movie}), the \metric{} metric might be helpful to tease apart the most favorable model. Our results indicate that \metric{} complements the standard sentence metrics in capturing human preference.

%
%
%
%
%

Additionally, attention lowers hallucination, but it appears that much of the gain from attention models is due to access to the underlying convolutional features as opposed the attention mechanism itself.
Furthermore, we see that models with stronger \emph{image consistency} frequently hallucinate fewer objects, suggesting that strong visual processing is important for avoiding hallucination.

Based on our results, we argue that the design and training of captioning models should be guided not only by cross-entropy loss or standard sentence metrics, but also by image relevance.
Our \metric{} metric gives a way to evaluate the phenomenon of hallucination, but other image relevance metrics e.g. those that incorporate missed salient objects, should also be investigated. We believe that incorporating visual information in the form of ground truth objects in a scene (as opposed to only reference captions) helps us better understand the performance of captioning models.

%
%

%

%

